\documentclass[instructions.tex]{subfiles}
\begin{document}
 
\chapter{De la durée des peines de l’enfer,}
 
Dieu ne punit jamais le péché en cette vie dans toute sa rigueur, parce que cette vie est le temps de sa miséricorde; il le punit en l’autre dans toute sa sévérité, parce que c’est le temps de sa justice. Il souffre dans le temps les crimes des hommes avec patience, dit Tertullien, parce qu’il aura l’éternité tout entière pour les punir\footnote{Patiens, qui æternus. Tert.}. S’il récompense dans le ciel, il récompense en Dieu; s’il punit en enfer, il doit punir en Dieu. Le bonheur des saints est souverain, le malheur des réprouvés l’est de même; leurs peines sont souveraines, parce qu’elles sont sans fin, sans consolation, sans adoucissement.
 
\primo Les peines de l’enfer sont sans fin, puisque l’éternité n’a point de terme dans sa durée. Tout ce qui s’en écoule ne la diminue point, elle est toujours tout entière après des millions d’années, toujours elle commence et ne finit jamais; ôtez de l’éternité autant de siècles qu’il y a de feuilles dans la forêt et d’atomes dans l’air, vous serez encore au commencement de l’éternité. Il y a plus de cinq mille ans que Caïn souffre dans l’enfer; il n’est pas plus avancé dans son éternité que le premier jour.
 
Quand un rocher seroit d’une masse égale au globe de la terre, et qu’un oiseau n’en transporteroit qu’un grain chaque siècle, ce rocher enfin changeroit de place: et après que cent mille rochers auroient été ainsi transportés, l’éternité seroit toujours la même. Elle ne diminue ni dans sa durée, ni dans ses rigueurs; le feu qui brûle les réprouvés seroit toujours aussi ardent, leurs douleurs aussi insupportables, leurs regrets et leurs larmes aussi amers, le ver de la conscience qui les déchire toujours aussi cruel\footnote{Vermis eorum non moritur, et ignis non extinguatur. Marc. 9. 43.}. Où est votre raison, si vous ne réfléchissez pas à cette vérité? et si vous y réfléchissez, dans quelle agitation ne doit pas être votre âme? Oh! que vous êtes profondément endormi dans les ombres de la mort, si cette vérité frappante ne vous réveille pas!
 
Si un juge vous menaçoit de faire lier votre main dans le feu pendant une heure, que ne feriez-vous pas pour éviter ce supplice? Dieu vous menace d’un feu qui ne s’éteindra jamais, et vous vivez sans précaution! vous craindriez le tourment d’une heure, quelle est votre stupidité de ne pas craindre un tourment éternel!
 
Quand on ne souffriroit en enfer qu’un léger mal de tête toute l’éternité, on devroit tout sacrifier pour s’en préserver; disons plus, il n’est point d’homme qui voulût accepter un royaume à condition d’être, pendant trente ans, dans une prison. Où est donc votre bon sens de vous exposer à être dans un étang de feu et de soufre, non pas trente ans, mais pendant des années sans fin; non pas pour un royaume, mais pour un plaisir brutal, pour une débauche, pour un procès de chicane, pour quelques désirs de vengeance, etc.?
 
Pensez-en tout ce que vous voudrez; vous êtes dans un état bien déplorable, si vous ne changez de conduite.
 
\secundo Comprenez une bonne fois qu’on souffre en enfer des tourmens qui ne finissent jamais, et qu’on les souffre sans consolation. Le passé et l’avenir accablent un réprouvé; il est continuellement dans des sentimens de rage proportionnés à l’éternité de son malheureux sort; il ne peut s’empêcher de réfléchir à tout moment sur sa vie passée, et sur cette durée de siècles qui ne passeront jamais; il compare sans cesse, dans le repentir le plus affligeant, ce qu’il a été avec ce qu’il est, ce qu’il a fait avec ce qu’il souffre, ses plaisirs passés avec ses tourmens, sa vie si courte avec son éternité. Après avoir souffert un million de siècles, sa vie ne lui paroîtra qu’un moment, et à peine se souviendra-t-il qu’il a été au monde; il ne s’en souviendra que pour gémir sur ce qu’il a perdu, sur ce qu’il a fait et sur ce qu’il a mérité: Je suis damné, je le suis pour jamais; j’ai perdu mon Dieu, je l’ai perdu pour toujours! regrets accablans, regrets éternels, mais regrets infructueux!
 
\tertio Le réprouvé souffre de la sorte sans adoucissement, sans relâche. Il y a plus de 1800 ans que le mauvais riche pousse des cris lamentables, sans que ses larmes et ses gémissemens aient pu calmer un seul moment ses douleurs, ni lui obtenir une goutte d’eau.
 
O état déplorable de ne pouvoir, depuis tant d’années, obtenir si peu de chose par tant de cris et de pleurs! Que pense-t-il à présent des biens qu’il a possédés, des plaisirs dont il a joui, des injustices et des crimes qu’il a commis? Qu’en penserez-vous vous-même un jour? Si les rochers étoient capables de raison et de sentiment, ils deviendroient sensibles à ces vérités. Concluons que quiconque ne s’applique pas à se sauver pendant que Dieu lui en donne le temps, est bien aveugle et bien stupide.
 
O vous, qui lisez ou entendez ces vérités! souvenez-vous que les jours qui vous restent sont courts. Vous n’avez point de temps à perdre; si vous le négligez, arrivera un moment auquel vous ne pourrez plus y penser. Vous pouvez à présent vous convertir à Dieu et vous préserver des feux éternels; mais si vous êtes un jour condamné, l’enfer sera la maison de votre éternité. On ne se perd qu’une fois, et c’est pour toujours\footnote{Periisse semel æternum est.}. Qui est perdu, l’est pour jamais. Mon Dieu! y pensons-nous?
 
\section{Résolutions}
 
\primo Puisque les rigueurs que la justice divine exerce sur les réprouvés, sont si terribles, et qu’elles ne finiront jamais, il faut qu’à quelque prix que ce soit, je les évite.
 
\secundo En conséquence, je renonce dès ce moment à ces fréquentations, à ces lectures, à ces attaches dangereuses.
 
\tertio Je remplirai aussi et avec toute la fidélité dont je suis capable, ces devoirs qui répugnent tant à mon orgueil, à ma lâcheté, à ma trop grande délicatesse.
 
\quarto Et je descendrai souvent en enfer en esprit, afin de ne pas y descendre en effet après ma mort. 

\prayer{Mon Dieu, ayez pitié de moi, et ne permettez pas que je me perde pour l’éternité.}
 
\end{document}
